% !TEX TS-program = xelatex
% !TEX encoding = UTF-8 Unicode
% !Mode:: "TeX:UTF-8"

\documentclass{resume}
\usepackage{zh_CN-Adobefonts_external} % Simplified Chinese Support using external fonts (./fonts/zh_CN-Adobe/)
% \usepackage{NotoSansSC_external}
% \usepackage{NotoSerifCJKsc_external}
% \usepackage{zh_CN-Adobefonts_internal} % Simplified Chinese Support using system fonts
\usepackage{linespacing_fix} % disable extra space before next section
\usepackage{cite}
\usepackage{bookmark} % fix outline/bookmark warnings, load after hyperref (from resume.cls)

\begin{document}
\pagenumbering{gobble} % suppress displaying page number

\name{江裕诚}

\basicInfo{
  \email{fatjyc@gmail.com} \textperiodcentered\ 
  \phone{(+86) 186-666-26214}}

\section{\texorpdfstring{\faStar\ 个人亮点}{个人亮点}}
\begin{itemize}[parsep=0.5ex]
  \item 10+ 年平台与开发者工具研发经验，长期负责代码托管、版本管理、工作流、私有化交付与复杂系统架构演进。
  \item 具备 Agent、RAG、工作流编排与执行验证闭环的工程落地经验，参与 AI 测试开发 Agent 与 AI 开发平台建设。
  \item 擅长将模型能力转化为可执行、可控、可观测的系统能力，关注 Tool Use、State、Workflow、Failure Handling 与平台集成。
  \item 长期在真实开发场景中使用 Cursor、Cloud Code、Codex 等 AI 开发工具，覆盖代码生成、重构、调试与问题排查，具备将 AI 工具稳定融入研发工作流的实践经验。
  \item 跨栈研发能力强（Java/TypeScript/Ruby/Rust、React、IDE 插件/在线开发环境/CLI），并具备 B2B 解决方案与企业级交付经验，能够从底层基础设施到上层 AI 应用做跨层设计与实现。
\end{itemize}

\section{\texorpdfstring{\faCogs\ 技能}{技能}}
% increase linespacing [parsep=0.5ex]
\begin{itemize}[parsep=0.5ex]
  \item \textbf{编程语言/框架}: Java、JavaScript/TypeScript、Ruby、Rust、React
  \item \textbf{LLM/Agent 工程}: LLM 应用开发、Agent 集成、向量数据库（RAG）、工作流编排与调度
  \item \textbf{DevOps/平台}: Git、CI/CD、Docker、Kubernetes
\end{itemize}

\section{\texorpdfstring{\faBuilding\ 工作经历}{工作经历}}

\datedsubsection{\textbf{深圳市巴别科技有限公司}}{2022年5月 -- 2025年7月}
\begin{itemize}
  %% [MOD] 从“负责/参与”改为更清晰的职责边界；去掉分号；补足“平台/多产品线”表达
  \item \roletitle{开发工程师}{负责 AI 驱动的云原生开发平台与 Agent 相关能力研发；参与前后端服务设计与落地，支撑 gru.ai、babel.cloud、gbox.ai 等多产品线演进。}
\end{itemize}

\datedsubsection{\textbf{深圳市腾云扣钉科技有限公司}}{2014年4月 -- 2022年2月}
\begin{itemize}
  %% [MOD] 强化“负责人”影响范围；统一术语（代码评审/MR/Webhook）；提升表达专业度
  \item \roletitle{代码托管业务开发负责人}{主导代码托管业务线技术规划与团队协作机制建设；负责 Git 协议栈、代码评审（MR）、Webhook 等核心能力的架构演进；带领团队完成私有化部署、OEM 定制等企业级方案交付。}
  %% [MOD] 调整措辞，强调产品矩阵覆盖面
  \item \roletitle{平台开发工程师}{参与 CODING（https://coding.net/）多业务线研发，覆盖项目管理、代码托管、CI/CD 等 DevOps 产品矩阵。}
  %% [MOD] 统一“参与/负责”风格，避免口语化
  \item \roletitle{码市平台开发工程师}{参与码市（https://codemart.com/）核心功能研发，承担需求评审、技术方案设计与关键模块实现。}
  %% [MOD] “对接”更偏执行，补充“架构/交付闭环”更贴近企业方案岗位价值
  \item \roletitle{企业客户解决方案工程师}{负责企业客户售前技术咨询、POC 验证与私有化部署方案设计；推进部署交付落地与问题闭环。}
\end{itemize}

\datedsubsection{\textbf{深圳市奥思网络科技有限公司}}{2013年7月 -- 2014年3月}
\begin{itemize}
  %% [MOD] 精炼句式，突出模块范围
  \item \roletitle{开发工程师}{参与 Gitee（https://gitee.com/）平台研发，负责代码托管、Issues、Wiki 等模块迭代与性能优化。}
\end{itemize}

\section{\texorpdfstring{\faUsers\ 项目经历}{项目经历}}

\datedsubsection{\textbf{gru.ai}}{2023年10月 -- 2025年7月}
https://gru.ai \ | \ 面向 GitHub Pull Request 流程的 AI 测试开发平台
\role{开发工程师}{}
\begin{onehalfspacing}
\begin{itemize}
  \item 参与构建面向 GitHub Pull Request 流程的执行型 AI 单元测试开发 Agent，将代码变更理解、上下文检索、测试生成、sandbox 验证与 PR 回写串联为完整闭环
  \item 基于统一 Agent Behavior 框架实现配置驱动的测试 Agent 流程，抽象 job / plan / task / step 执行结构，支持多步任务编排、失败恢复与终止控制
  \item 接入 GitHub API、代码搜索与依赖解析能力，补充目标文件调用关系、相关源码与外部依赖上下文，减少因上下文不足导致的生成失败与无效重试
  \item 构建基于容器 / sandbox 的执行验证环境，支持 Node.js（npm）、Java（Maven）、Rust（cargo test）等技术栈的标准测试体系运行，以真实执行结果而非文本生成结果作为成功标准
  \item 设计重试阈值、自动终止和流程回写机制，推动 AI 从代码建议工具演进为面向研发流程的执行型 Agent 能力
\end{itemize}
\end{onehalfspacing}

\datedsubsection{\textbf{babel.cloud}}{2022年5月 -- 2025年7月}
https://babel.cloud \ | \ 面向 AI 场景的云原生开发平台
\role{开发工程师}{}
\begin{onehalfspacing}
\begin{itemize}
  \item 参与构建面向 AI 场景的云原生在线开发平台 Babel Cloud，负责平台组件编排、代码版本管理及向量检索基础设施建设
  \item 设计并实现平台组件组装能力，将 HTTP、Function、Worker 等 TypeScript 模块编排为完整可运行的 Node.js 服务，并注入调试、监控等运行时能力
  \item 基于 Git 与 Rust 相关 Git 库实现平台内代码版本管理能力，支持代码存储、历史追踪、版本回退与分支操作，打通底层版本语义与上层工作区体验
  \item 基于 PostgreSQL 及向量插件实现平台内建向量数据库能力，并设计类 Pinecone 风格 API，支撑用户在平台内直接开发 AI 相关应用
  \item 推动平台从“在线开发环境”演进为“应用构建、代码管理与 AI 能力一体化”的云原生开发平台底座
\end{itemize}
\end{onehalfspacing}

\datedsubsection{\textbf{CODING 平台}}{2013年4月 -- 2022年2月}
https://coding.net \ | \ 一站式 DevOps 平台（代码托管、CI/CD、项目管理）
\role{开发工程师，代码托管业务开发负责人}{}
\begin{onehalfspacing}
\begin{itemize}
  %% [MOD] 统一术语；“支撑”改为“提供/建设”更主动
  \item 负责代码托管业务的架构设计与研发，建设代码托管、代码评审（MR）、Webhook 等核心能力
  %% [MOD] “B2C 转 B2B”改为“企业化转型”更职业；拆分点更清晰
  \item 主导产品企业化转型：前端从 AngularJS 迁移至 React；后端按服务域拆分并改造 API 网关，支撑多租户与权限体系升级
  %% [MOD] 语义更精确：多租户数据隔离/配置中心/离线方案属于交付关键项
  \item 主导 OEM 与私有化交付改造：设计多租户数据隔离、配置中心与离线部署方案，支撑多家大型企业客户交付
  %% [MOD] “MR、Webhook、大文件存储”统一风格；“推进团队建设”具体化为“流程与评审机制”
  \item 担任代码托管负责人：推进 Git 协议栈、MR、Webhook、大文件存储等能力演进；建立技术评审与迭代规划机制，提升交付效率与质量
\end{itemize}
\end{onehalfspacing}

\datedsubsection{\textbf{nocalhost}}{2020年9月 -- 2021年1月}
https://nocalhost.dev \ | \ 云原生远程开发工具（CNCF Sandbox）
\role{开发工程师}{}
\begin{onehalfspacing}
\begin{itemize}
  %% [MOD] “主导”与交付物绑定；“系列 IDE 插件”更精确
  \item 主导前端模块与 JetBrains 系列 IDE 插件（IntelliJ IDEA、GoLand 等）研发
  %% [MOD] 句式更聚焦“能力 + 价值”
  \item 实现端口转发、文件同步、远程调试、热加载等能力，使开发者可在本地 IDE 中透明操作集群内工作负载
\end{itemize}
\end{onehalfspacing}

\datedsubsection{\textbf{码市平台}}{2015年1月 -- 2016年6月}
https://codemart.com \ | \ 互联网软件外包服务平台（交易撮合、支付、评价）
\role{开发工程师}{}
\begin{onehalfspacing}
\begin{itemize}
  %% [MOD] 更工程化：列举“关键模块”并保持并列结构
  \item 参与平台建设与迭代：负责账号体系、发布/接单流程、支付系统（支付宝/微信）、管理后台等核心模块
\end{itemize}
\end{onehalfspacing}

\datedsubsection{\textbf{Gitee}}{2013年7月 -- 2014年3月}
https://gitee.com \ | \ 国内代码托管平台（Git 托管、Issues、Wiki）
\role{开发工程师}{}
\begin{onehalfspacing}
\begin{itemize}
  %% [MOD] 与“工作经历”避免重复，但保留项目维度描述；更简洁
  \item 参与代码托管、Issues、Wiki 等模块研发与性能优化，支持平台功能迭代与稳定性提升
\end{itemize}
\end{onehalfspacing}

\section{\faGraduationCap\ 教育背景}
\datedsubsection{\textbf{桂林电子科技大学信息科技学院}}{2009 -- 2013}
\textit{工学学士}\ 计算机科学与技术 \textperiodcentered\ 广西桂林

\end{document}
